%! Author = Marjan
%! Date = 06/07/2025

% Preamble
\documentclass[11pt]{book}

% Packages
\usepackage{amsmath}    % For math symbols and equations
\usepackage{graphicx}   % For including graphics
\usepackage{geometry}   % For adjusting page layout
\usepackage{fancyhdr}   % For custom headers/footers
\usepackage{hyperref}   % For hyperlinks
\usepackage{listings}   % For code listing
\usepackage{lipsum}     % For placeholder text (to test layout)
\usepackage{helvet}
\usepackage{arydshln}
\usepackage{tikz}        % For creating diagrams and drawing
\usetikzlibrary{
    positioning,
    arrows.meta,
    shapes,
    fit,
    backgrounds,
    calc
}

\usepackage{pgfplots}    % For plotting graphs
\usepackage{wasysym}
\usepackage{float}
\usepackage{coptic}
\usepackage{xcolor}
\usepackage{tweaklist}
\usepackage{tikzexternal}
\usepackage{listings}
\usepackage{xcolor}

\definecolor{terminalbg}{RGB}{30,30,30}
\definecolor{terminalfg}{RGB}{220,220,220}
\definecolor{terminalgreen}{RGB}{152,195,121}

\lstdefinestyle{terminal}{
    backgroundcolor=\color{terminalbg},
    basicstyle=\ttfamily\small\color{terminalfg},
    frame=single,
    rulecolor=\color{black},
    breaklines=true,
    columns=fullflexible,
    showstringspaces=false,
    keywordstyle=\color{terminalgreen},
    commentstyle=\color{gray},
}


% Dockerfile language definition
\lstdefinelanguage{Dockerfile}{
    sensitive=true,
    keywords={
        FROM,RUN,CMD,LABEL,MAINTAINER,EXPOSE,ENV,ADD,COPY,
        ENTRYPOINT,VOLUME,USER,WORKDIR,ARG,ONBUILD,STOPSIGNAL,
        HEALTHCHECK,SHELL
    },
    keywordstyle=\color{blue}\bfseries,
    comment=[l]{\#},
    commentstyle=\color{green!50!black}\itshape,
    stringstyle=\color{red},
    morestring=[b]",
    morestring=[b]',
}

% Style for Dockerfiles
\lstdefinestyle{dockerstyle}{
    language=Dockerfile,
    basicstyle=\ttfamily\small,
    keywordstyle=\color{blue}\bfseries,
    commentstyle=\color{green!50!black}\itshape,
    stringstyle=\color{red},
    numbers=left,
    numberstyle=\tiny\color{gray},
    stepnumber=1,
    numbersep=8pt,
    frame=single,
    breaklines=true,
    breakatwhitespace=true,
    showstringspaces=false,
    tabsize=2,
    captionpos=b,
    columns=fullflexible
}


%\usepackage{amstex} % Helvetica as a sans-serif font alternative
\renewcommand{\rmdefault}{phv} % Set the default font family to Helvetica

% Set page margins (A4 paper)
\geometry{top=1in, bottom=1in, left=1in, right=1in}

% Set up the header
\pagestyle{fancy}
\fancyhf{}
\fancyhead[L]{Docker Notes}
\fancyhead[C]{Your Name}
\fancyhead[R]{\thepage}
% Document
\begin{document}
% Title Page
    \begin{titlepage}
        \centering
        \vspace*{2in}
        \Huge \textbf{Docker Notes}
        \vfill
        \Large Your Name
        \vfill
        \Large Date: %\today
    \end{titlepage}

    \newpage

    \tableofcontents
    \newpage


    \chapter{Notes info introduction}

    \paragraph{Purpose of notes}
    These are a set of notes on Docker.


    \chapter{Geting Started with Docker}


    \section{Base CLI Command}

    \subsection{Docker Image commands}
    \begin{lstlisting}[style=terminal]
        docker image build
        docker image history
        docker image inspect
        docker image pull
    \end{lstlisting}

    \subsection{Docker Container Commands}
    \begin{lstlisting}[style=terminal]
        docker container create
        docker container exec
        docker container logs
        docker container run
    \end{lstlisting}

    Can use man pages via CLI.
    See \texttt{docker container --help}
    See \texttt{docker container create --help}

    Consider old syntax vs new syntax
    \begin{lstlisting}[style=terminal]
        docker run <image>
        docker ps
        docker logs <ps>
        #--------------New syntex
        docker container run
        docker container ls
        docker contiainer logs <id>
    \end{lstlisting}

    \texttt{docker ps} is the same as \texttt{docker container ls}


    \section{Images vs Containers}
    A Docker image is a standalone, read-only package that includes everything needed to run a software application: the code, runtime, system tools, libraries and settings.
    \begin{itemize}
        \item App Runtime
        \item Libraries
        \item Config and Env Vars
        \item Dependencies
    \end{itemize}

    Docker containers are running instances of images.
    See docker hub for more

    Can launch multiple containers from a Single Docker Image
    Can launch multiple containers from different docker Images

    \subsection{Attached vs Detached modes}
    In most cases you will be running containers in detached modes, however there some specific use cases where you might wan tot run it in attached mode as well.

    When you use the docker container run <image> without any extra flags, Docker starts the container in Attached Mode.
    In this mode, your terminal's local standard input, output and error streams are "attached" to the container.

    Note: If you press Ctrl + C in ther terminal, you send a SIGINT signal to the container's primary process, which usually stops the container entirely.

    To run a container in the background, you use the -d or --detach flag.
    This detached mode is the standard for any long-running service, for example web server containers or database containers.
    Ideally, you would like it to be runnig all of the time so you would use it with detached mode.


    \begin{lstlisting}[style=terminal]
        docker container run nginix
        docker container run --detach nginix
        docker container run -d nginix
    \end{lstlisting}

    Generally in most cases, you wil be using the detached mode in the production environments.

%    TODO add detached mode and attached mode table for reference

    \subsection{Basic Container Management Commands}

    When you have container running you can perform multiple operations on the container.

    \subsubsection{Managing the state of a container}
    \begin{lstlisting}[style=terminal]
        docker container start <name>
        docker container stop <name>
        docker container restart <name>
    \end{lstlisting}

    \subsubsection{Monitoring Performance}
    Example: A container can be consuming a certain level of memory and CPU
    \begin{lstlisting}[style=terminal]
        docker container stats <name>
    \end{lstlisting}

    This command can be useful in identifying which container is using the most CPU.

    \begin{lstlisting}[style=terminal]
        docker container inspect <name>
    \end{lstlisting}

    This command will get detailed information related to the container.

    \begin{lstlisting}[style=terminal]
        docker container rm <name>
        docker container rm <name> --force
        #--- use for exited containers
        docker container ls -a
        docker container rm <name1> <name2>
    \end{lstlisting}

    This will remove a container. If the container is running you can use force remove.

%    TODO add table

    \subsection{Basic Container Lifecycle}
    A docker container lives as long as its main process runs.
    If the process stops, the container stops.
    If the process is running the container will continue to run.

    Short-lived containers
    - Use for scripts, batch jobs and one time commands.

    Long Running Containers
    - Run a continuous process
    User for Web Servers, APIs, Database etc.

    In production environments you will most likely use long-running containers.

    \subsubsection{Default Container Command}
    In Docker, the default container command is the instruction that tells the container process which process to run immediately upon starting.
    You can override the default container command if required.

    \subsection{Docker Container exec}
    \texttt{docker exec} lets you run a command inside an already running container.

    \begin{lstlisting}[style=terminal]
        docker container exec nginx ls -l
    \end{lstlisting}

    \subsection{docker container exec - important of -i and -t flags}

    For simple commands that just print the output, you don't need any flags. The command runs, prints its result to your terminal and exits.
    No interaction is needed.

    In most cases however you would want ot go inside the container, so you will require a shell.
    In such cases, you will require an additional set of flags that need to be added here.

    If you want to connect inside a container using a shell like bash, sh (similar to SSH), the default approach will not work.
    Bash will close immediately.
    This is because it is an interactive shell that detects there is no input stream (STDIN) to read from.
    Without an active input stream or a terminal, it assumes its job is done and exists immediately.

    \subsection{The -i flag (interactive)}
    By default, Docker runs commands in non-interactive mode and does not keep the input stream (STDIN) open.
    using -i keeps STDIN attached to the container process, which is necessary for any interactive process.

    \begin{lstlisting}[style=terminal]
        docker container exec -i nginx bash
    \end{lstlisting}

    \subsection{The -t Flag (TTY)}
    This will give you a full-fledged experience.
    Overview: the hyphen t flag allocates a pseudo terminal.
    If you just use the -t flag, the standard input get closed so that the keystrokes have no path to reach the process inside the container.

    \begin{lstlisting}[style=terminal]
        docker container exec -t nginx bash # This is very similar to SSH into a Linux server
    \end{lstlisting}

    Using the -it option, users can get a full interactive shell.
    You can type, navigate and explore.

    \begin{lstlisting}[style=terminal]
        docker container exec -it nginx bash
        docker container exec -i -t nginx bash #Equivalent
        exit # Return to the host system.
    \end{lstlisting}

    Using this command, you can run everything inside the container environment.

    Note: Not every container comes with a shell, however the majority of them do.
    Note: Navigate Linux/OS file directory to see which binaries are available.


    \section{Basics of Ports}
    When data arrives at your computer or a server, the operating system needs to know which program should receive it.
    Ports help figure out which application should receive incoming data.
    An application can listen on a specific port number (port binding)

    When incoming data comes, this data will also have the port number.
    The server will forward the data to the application with the matching port number.
    Ports are numbers from 0 to 65,535 and only one process can bind to a specific port on the same IP address at a time.

    It is the responsibility of the incoming data to provide the port number so that the operating system can correctly map and send the data to the appropriate program over here.


    \begin{lstlisting}[style=terminal]
        netstat -ntlp # for linux. See what applications are bound to which ports.
    \end{lstlisting}


    \section{Publishing Ports in Docker}
%    TODO mebe add a diagram
    When a container has been started, it is assigned a unique IP address within its configured Docker network.
    The application running inside the container may listen on a specific port.

    Containers connected to the same network can communicate with each other using the contain's IP address and the application's listening port.

    Note: The internal Docker network IP addresses are not directly accessible from the host system (In Windows and Mac OS).
    On linux hosts, the Docker bridge network is directly reachable from the host via container IP.

    \begin{lstlisting}[style=terminal]
        curl 172.17.0.2:80 # This will work on Linux not on windows or mac OS
    \end{lstlisting}

    \subsection{Ports Concept}
    When we talk about ports in Docker, we are always talking about two different locations.
    - Host Port (outside) - The port you type into your browser (e.g. localhost:8080)
    - Container Port (inside) - The port the actual application is configured to use (e.g. 80)

    \subsection{Connecting to Container Application}
    To allow a container to accept incoming connections from the host or the outside world, you must bind (map) it to a host port.
    Once this is done the users can send a request to the IP address of the server followed by the port number.

    \begin{lstlisting}[style=terminal]
        docker container run -d --name nginx -p 8080:80 nginx
        curl 172.17.0.2:8080 # this should work after the above command
    \end{lstlisting}

    \section{Docker Restart Policy}
    by default, Docker containers will no start when they exit or when the docker daemon is restarted.
    Sometimes the server restarts and we want all the the containers to automatically be started again.

    Docker provides restart policies to control whether your containers start automatically when they exit, or when Docker restarts.

    \subsection{Restart policies}
    You can specify the restart policy by using the --restart flag with the docker run command.
    \begin{itemize}
        \item no - The defaul. Docker will not restart the container under any circumstances
        \item on-failure - Restarts only if the container exits with a non-zero exit code
        \item unless-stopped - Always restarts except when you manually stop the container. If manually stopped, it won't restart even after a daemon restart.
        \item always - Always restarts regardless of exit code, including after a daemon restart --- even if the container way manually stopped before the daemon went down.
    \end{itemize}

    In production environments we want the containers to start back up again should the server restart.

    \subsection{Choosing the right policy}
    \begin{table}[ht]
        \centering
        \begin{tabular}{|c|c|c|c|c|}
            \hline
            Policy         & Restarts on crash? & Restarts on success exit? & Restarts after daemon restart & Best for \\ \hline
            no             & No                 & No                        & No                            & Development, testing           \\ \hline
            always         & Yes                & Yes                       & Yes                           & Critical services, web servers \\ \hline
            on-failure     & yes                & no                        & Only if it was running        & Batch jobs, scripts            \\ \hline
            unless-stopped & Yes                & Yes                       & Only if no manually stopped   & Databases, background services \\ \hline
        \end{tabular}
%        \caption{Your table caption}
%        \label{tab:mytable}
    \end{table}


    \chapter{Creation, Management and Registry}


    \section{Basics of a Dockerfile}
    A docker container is a running instance of image, whatever container you make is based of one of these docker images.
    There are numerous images available. From these images, multiple containers can be created.

    For an organisation, you may want your own customized images.

    A dockerfile is a simple text file containing ordered step-by-step instructions (commands) used to build a Docker image.

    The workflow is, docker file into docker image into a running docker container.
    To create a docker file create a a file called "Dockerfile".

    \begin{lstlisting}[style=terminal]
        docker build -t image_name . # will reference the docker file in the same folder
        docker container run -d -p 8080:80 --name container_name container_image:latest #Create the container
    \end{lstlisting}

    There is a wide range of instructions available that can be used in a Dockerfile to build and image.

    Some of these include:

    \begin{table}[ht]
        \centering
        \begin{tabular}{|c|c|}
            \hline
            Instructions & Description                                                            \\ \hline
            FROM         & The starting point. Every Dockerfile must start with this              \\ \hline
            CMD          & The default command that executes                                      \\ \hline
            RUN          & Executes a command during the build process (e.g installing a package) \\ \hline
            COPY         & Copies files from your local machine into the container                \\ \hline
            SHELL        & Set the default shell of an image                                      \\ \hline
        \end{tabular}
%        \caption{Your table caption}
%        \label{tab:mytable}
    \end{table}

    \subsection{Format of a Dockerfile}
    The format of a Dockerfile follows the syntax INSTRUCTION <arguments>


    \section{Docker file Instructions}
    Before you can createa a docker file the necessary set of instructions.
    See Docs for the set of instructions. Using this you can build a custom set of images.

    \subsection{FROM}
    The FROM instruction sets the base image that all subsequent instructions build upon.

    \subsection{RUN}
    The RUN instruction executes a command during the build process
    \texttt{RUN apt-get install nginx}

    \subsection{CMD}

    CMD instruction allows users to provide the default command that executes when the container starts.

    \subsection{Example DockerFile}
    Prior to building a docker image, try it in a base container.

    \begin{lstlisting}[style=terminal]
        docker container run -d --name ubuntu-container ubuntu:latest sleep 3600 # Run the container
        docker container exec -it ubuntu-container bash # Open a bash shell on the container to check if desired programs are installed

        # --------- from inside the container -----
        apt update
        apt install iputils-ping -y # Install ping for this example. -y so no prompt is asked.
    \end{lstlisting}

    Now we can map this to instructions in the DockerFile.

    \begin{lstlisting}[style=dockerstyle,caption={Example Dockerfile}]
        FROM ubuntu:latest

        RUN apt update

        RUN apt install iputils-ping -y

        CMD ["ping", "google.com"]
    \end{lstlisting}

    Using this image file we can build an image.

    \begin{lstlisting}[style=terminal]
        docker build -t <image-name> . # Ensure that you are in the right folder
        docker container run --name <container-name> <container-image>
    \end{lstlisting}

    \subsection{COPY and ADD}
    COPY and ADD serve the same fundamental purpose
    "they can move files from you host machine into the container's file system"

    \begin{table}[ht]
        \centering
        \begin{tabular}{|c|c|c|}
            \hline
            Feature         & COPY & ADD \\ \hline
            Primary Function             & Copies local files / folder to the image                 & Copies local files / folders to the image    \\ \hline
            Remote URLs         & Not Supported                & Can download files from a URL   \\ \hline
            Auto-Extraction     & Does not extract compressed file                & Automatically extracts .tar, gzip, etc    \\ \hline
            Best Practice & Recommended from most use cases & Use only when specific features are needed   \\ \hline
        \end{tabular}
%        \caption{Your table caption}
%        \label{tab:mytable}
    \end{table}

    Think of COPY as a simple "Ctrl+C, Ctrl+V.". It is transparent and predictable.
    Docker official documentation recommends using COPY unless you have a specific reason not to.

    \subsection{ENTRYPOINT}
    \subsubsection{Revising Basics of CMD Instruction}
    CMD allows users to provide the default command that executes when the container starts.
    CMD instructions can easily be over run by the docker run arguments.

    \subsubsection{ENTRYPOINT instruction}
    ENTRYPOINT sets the default executable for your container.

    Any arguments supplied to the docker run command are appended to the ENTRYPOINT command

    \begin{lstlisting}[style=dockerstyle,caption={Example Dockerfile}]
        FROM ubuntu:latest

        RUN apt update

        RUN apt install iputils-ping -y

        ENTRYPOINT ["ping"]
    \end{lstlisting}

    \begin{lstlisting}[style=terminal]
        docker container run --name ping-container-1 ping-image google.com # google.com is appended to the argument list of ENTRYPOINT
    \end{lstlisting}

    Notes:
    Use ENTRYPOINT when you need your container to always run the same base command and you want to allow users to append additional commands at the end.
    ENTRYPOINT can be overridden on the docker rtun command line by supplying the \texttt{--entrypoint} flag.

    \begin{table}[ht]
        \centering
        \begin{tabular}{|c|c|c|}
            \hline
            Command         & Description & Use-case \\ \hline
            CMD             & Defines the default executable of a Docker image. It can be overridden by argumetns & Utility images allow users to pass different executables and arguments on the command line    \\ \hline
            ENTRYPOINT      & Defines the default executable. It can be overriden by the "--entrypoint" docker run arguments & Images built for a specific purpose where overriding the default executable is desired.    \\ \hline
        \end{tabular}
%        \caption{Your table caption}
%        \label{tab:mytable}
    \end{table}

    \subsection{WORKDIR}
    WORKDIR sets the working directory inside the container for any subsequent instructions (RUN, CMD, ENTRYPOINT, COPY, ADD)
    Without WORKDIR, instructions like RUN, CMD, ENTRYPOINT, COPY and ADD execute from the container's root directory (/)

    \begin{lstlisting}[style=dockerstyle,caption={Example Dockerfile}]
        FROM ubuntu:latest

        COPY config.json /my-application #Specifying the folder is bad practice

        ADD http://kplabs.in/friend.jpg /my-application

        RUN cd /my-application && npm install

        RUN cd /my-application && node server.js
    \end{lstlisting}

    \begin{lstlisting}[style=dockerstyle,caption={Example Dockerfile}]
        FROM ubuntu:latest

        WORKDIR /my-application # Good practice

        COPY config.json

        ADD http://kplabs.in/friend.jpg

        RUN npm install

        CMD ["node", "server.js"]
    \end{lstlisting}

    The WORKDIR instruction can be used multiple times in a Dockerfile.

    \subsection{EXPOSE}
    To connect to a container from the outside network, you must publish its port by mapping it to a port on the host machine (e.g. docker run -p 8080:80)

    EXPOSE act as a documentation label that says "This container listens on this port"
    It does not actually open any port to the outside world.

    Note:
    If you use the -P flag when running a container (docker run -P), Docker will automatically map all EXPOSED ports to random high-order ports on your host machine.

    \subsubsection{Useful command}
    \begin{lstlisting}[style=terminal]
        docker image inspect nginx-image
    \end{lstlisting}
    
    \subsection{ARG}
    ARG instuction allows us to set build-time variables.

    It allows you to pass specific values for a variable while it is being built.

    \begin{lstlisting}[style=dockerstyle,caption={Example Dockerfile}]
        FROM ubuntu:latest

        ARG CONFIG_VERSION=1.2.0

        RUN touch tmp-v${CONFIG_VERSION}.json /tmp

        CMD ["sleep", "3600"]
    \end{lstlisting}

    \begin{lstlisting}[style=terminal]
        docker build --build-arg CONFIG_VERSION=1.2.5 -t test-image-02 . # Overrride the argument in the command line
    \end{lstlisting}

    The values of ARG do not persist once the image is built and the container is running.

    \subsection{ENV}

    The ENV instruction sets the environment variable inside a Docker image.
    These values will persist.

    \begin{lstlisting}[style=dockerstyle,caption={Example Dockerfile}]
        FROM ubuntu:latest

        ENV PLANET="world"

        CMD echo "Hello $PLANET"
    \end{lstlisting}

    \subsubsection{Setting and Overriding the Variables}
    You can use the -e, --env and --env-file flags to set simple environment variables in the container you're running,
    or overwrite variables that are defined in the Dockerfile of the image you're running.

    \begin{lstlisting}[style=terminal]
        docker container run --env PLANET=Pluto env-image-01
    \end{lstlisting}

    \subsection{LABEL}
    LABEL adds key-value metadata to a Docker image.

    Think of it like tags or sticky notes attached to your image --- they don't affect how the container runs but provide useful information for humans and tooling.

    \begin{lstlisting}[style=dockerstyle,caption={Example Dockerfile}]
        FROM nginx:alpine

        LABEL maintainer="maintainer maintainer@maintain.com"

        LABEL use="prototype"

        RUN echo '<h1>Hello</h1>' > /usr/share/nginx/html/index.html

        CMD ["nginx", "-g", "daemon off;"]
    \end{lstlisting}

    If you have hundreds of images on your server, you cna filter them by their labels.

    \begin{lstlisting}[style=terminal]
        docker images --filter "label=use=prototype"
    \end{lstlisting}
    
    \section{Docker Image Prune}
    Unused docker images can quietly take up valuable disk space, leading to slower system perofrmance and potential storage issues.
    Regularly pruning them helps maintain a clean, efficient and prodiction ready environment.


    \begin{lstlisting}[style=terminal]
        docker image prune # Only removes dangling images - A dangling image is one that isn't tagged and isn't referenced by any container.
        docker image prune # All imags not used by any container.
    \end{lstlisting}

    \section{Layers of a Docker Image}
    A Docker image is not a single large file.

    It's built as a stack of read-only layers, where each layer represents a discrete change.

    Note:
    Key instructions like RUN, COPY and ADD create layers in your final image.

    Once an image layer is built, it is read-only. It cannot be changed.

    \begin{lstlisting}[style=dockerstyle,caption={Example Dockerfile}]
        FROM ubuntu:latest

        RUN apt update

        RUN apt install nginx -y

        RUN apt install nano -y

        RUN echo '<h1>Hello from container</h1>' > /usr/share/nginx/html/index.html

        CMD ["nginx", "-g", "daemon off;"]
    \end{lstlisting}

    \begin{lstlisting}[style=terminal]
        docker image history <image-name>
    \end{lstlisting}

    0B instruction are considered meta data.

    The writeable layer:
    When a container runs, Docker adds one thin writable layer on top (the container layer).
    All changes made during the container's lifetime are stored in this thin R/W layer, leaving the underlying image layers untouched and immutable.

    One Image and Multiple Containers:
    If multiple containers are launched from the same image, each container will have its own indpendent Read/Write layer.
    These containers will need to be removed if you wish to delete the image.

    \section{Docker Build Cache}
    Docker images are built layer by layer from a Dockerfile.

    Each instruction that modifies the filesystem --- such as RUN,COPY and add --- adds a new read-only layer to the image stack.

    Docker Build Cache reuses these layers from previous builds when the instruction (and its inputs) haven't changed.
    This can significantly improve the overall build time for the docker image and is more efficient.

    Note:
    The build cache follows a chain: if a lyer is modified (e.g. a file change in a COPY instruction), every subsequent layer is invalidated and forced to rebuild.


    \section{Dockerfile - HEALTHCHECK instruction}
    HEALTHCHECK instruction - Docker allows us to tell the plaform how to test that our application is healthy.
    When Docker starts a container, it monitors the process that the container runs. If the process ends, the container exits.

    That's just a basic check and does not necessarily tell the detail about the application.

    \begin{lstlisting}[style=dockerstyle,caption={Example Dockerfile}]
        FROM busybox

        HEALTHCHECK --interval=5s CMD ping -c 1 172.17.0.2 #Health check runs every 5 seconds.
    \end{lstlisting}

    \subsubsection{Parameters for HEALTHCHECK}
    We can specify certain options before the CMD operation, these include:

    HEALTHCHECK --interval=5s CMD ping -c 1 172.17.0.2
    --interval=DURATION (default: 30s)
    --timeout=DURATION (default: 30s)
    --start-period=DURATION(default: 0s)
    -retries=N(default: 3)
    
    \section{HEALTHCHECK Instruction Options}
    HEALTHCHECK instruction is combined with various options to get the desired results.


    --interval=DURATION (default: 30s)
    --timeout=DURATION (default: 30s)
    --start-period=DURATION(default: 0s)
    -retries=N(default: 3)

    If you do not specify these options while adding a health check instruction, then the associate default values will automatically be used.

    \subsubsection{Exist Status}
    The command's exit status indicates the health status of the container The possible values are:
    0: Success --- The container is healthy and ready for use
    1: Failure --- The container is not working correctly
    2: Reserved --- Do not use this exit code

    \subsubsection{Example use-case}
    Check every five minutes or so that a web-server is able to serve the site's main page within three seconds:

    \begin{lstlisting}[style=dockerstyle,caption={Example Dockerfile}]
        HEALTHCHECK --interval=5m --timeout=3s \ CMD curl -f http://localhost/ || exit 1
    \end{lstlisting}

    Health checks can also be performed using the docker health check command. SEE DOCS

    \begin{lstlisting}[style=terminal]
        docker run -dt --name tmp --health-cmd "curl -f http://localhost" --health-interval=5s busybox sh
    \end{lstlisting}

    \subsubsection{Overview of CURL}
    CURL is used with -f option to fail silently.
    This is mostly done to better enable scripts etc to better deal with failed attempts

%    TODO add links to health check docker documentation

    \section{Docker commit}
    Whenever you make changes inside the container, it can be useful to commit a container's file chagnes or settings into a new image.
    By default, the container being committed and its processes will be paused while the image is comitted.

    \begin{lstlisting}[style=terminal]
        docker container commit CONTAINER-ID myimage01 # Syntax
    \end{lstlisting}

    \subsubsection{Change Option while Commiting}
    The --change option will apply Dockerfile instructions to the image that is created.

    Supported Dockerfile instructions:
    - CMD| ENTRYPOINT | ENV | EXPOSE | LABEL | ONBUILD | USER | VOLUME | WORKDIR

    Whenever you make changes inside the container, it can be useful to commit a container's file change or settings into a new image.
    By default, the container being comitted and its processes will be paused while the image is comitted.

    \begin{lstlisting}[style=terminal]
        docker container commit CONTAINER-ID myimage01
    \end{lstlisting}

    \section{Flattening Docker Images}
    Generally, the more layers and image has the more the size of the image.
    Some of the image size goes from 5GB to 10GB.
    Flattening an image to a single layer can help reduce the overall size of the image.

    \begin{lstlisting}[style=terminal]
        docker image history <image> # Get image size
        docker export myubuntu > myubuntudemo.tar # Export file
        cat myubuntu.tar | docker import - myubuntu:latest
        docker images # Verify the image is there
        docker image history ubuntu # Should see that the image has only one layer and reduced size
    \end{lstlisting}

    This is referred to as, flattening the image.

    \section{Overview of Docker Registries}
    A registry is similar to a github.
    Generally whenever you code for system administrator, it might be some kind of confgiruation management scripts like Ansible etc.
    Whenever you write those scripts it is recommedned to put it in a centralized location where other people can access as well.
    The code will remain safe there because if you store it locally and the machine fails then your code will be lost.
    This is the reason you store it in a centralized location.

    The same is true for the centralized docker registry.

    \subsection{Understanding Docker Registry}
    A Registry is a stateless, highly scalable server side application that store and lets you distribute Docker images.
    Docker Hub is the simplest example.
    There are various types of registry available, which includes:

    \begin{itemize}
        \item Docker Registry
        \item Docker Trusted Registry
        \item Private Registry (AWS ECR)
        \item Docker Hub
    \end{itemize}

    \begin{lstlisting}[style=terminal]
        docker tag ubuntu:latest localhost:5000/myubuntu
        docker images # check if the image is there
        docker push localhost:5000/myubuntu # This will push to a docker registry
        #----------- pulling an image from the repository
        docker pull localhost:5000/myubuntu
        docker image # Verify that the image has been pulled from the registry.
    \end{lstlisting}

    \section{Pushing Images to a Central Repository}

    \begin{lstlisting}[style=terminal]
        docker login # Login into docker hub
        docker tag busybox 5066/mydemo@v1
        docker push 5066/mydemo
        #----- verify in docker hub
        docker pull 5066/mydemo:v1 # pull image from docker hub
        docker image # verify
        docker logout # log out of docker hub
    \end{lstlisting}

    The same goes for other private repositories like ECT.
    If you want to push certain images to ECR or if you want to pull certain images from ECR then you will have to login.
    Only after then will you be able to pull or push the images.

    Specifically this is related to the private repositories.


    \section{Applying Filter for Docker Images}
    Searching for docker images in the docker registry.

    SEE DOCS also

    \begin{lstlisting}[style=terminal]
        docker search nginx
        docker search nginx --limit 5 # Will list the top five starred repos
        docker search nginx --fitler "is-official=true"
    \end{lstlisting}

    \section{Moving Images Across Hosts}
    Overview of the Docker save command:
    The docker save command will save one or more images to a tar archive

    \begin{lstlisting}[style=terminal]
        docker save busybox > busybox.tar
    \end{lstlisting}

    Overview of Docker Load:
    The docker load command will load an image from a tar archive

    \begin{lstlisting}[style=terminal]
        docker load < busybox.tar
    \end{lstlisting}


    \chapter{Networking}


    \section{Overview of Docker Networking}


    \section{Understanding Bridge Networks}


    \section{Implementing User-Defined Bridge Networks}


    \section{Understanding Host Network}


    \section{Implementing None Network}


    \section{Publish All Argument for Exposed Ports}


    \section{Legacy Approach for Linking Containers}


    \chapter{Orchestration}
%    TODO consider skipping


    \chapter{Orchestration using Kubernetes}
%    TODO consider skipping


    \chapter{Installation and Configuration of Docker EE}
%    TODO consider skipping


    \chapter{Security}


    \section{Overview of Container Security Scanning}


    \section{Configuring Container Scan with DTR}


    \section{DTR Webhooks}


    \section{Overview of UCP Client Bundle}


    \section{Integrating CLI with UCP Client Bundle}


    \section{Overview of LDAP}


    \section{Integrating LDAP with UCP}


    \section{Linux Namespaces}


    \section{Control Groups (cgroups)}


    \section{Limiting CPUs for Containers}


    \section{Reservation vs Limits}


    \section{Swarm MTLS}


    \section{Managing Secrets in Docker Swarm}


    \section{Docker Content Trust}


    \section{Overview of Docker Groups}


    \section{Overview of Linux Capabilities for Docker}


    \section{Privileged Containers}


    \chapter{Storage and Volumes}


    \section{Overview of Docker Storage Drivers}


    \section{Block vs Object Storage}


    \section{Changing Storage Drivers}


    \section{Overview of Docker Volumes}


    \section{Bind Mounts}


    \section{Automatically Remove Volume on Container Exist}


    \section{Overview of Device Mapper}


    \section{Logging Drivers}


    \section{Creating Volumes in Kubernetes}


    \section{PersistentVolumes and PersistentVolumeClaims}


    \section{Volume Expansion in Kubernetes}


    \section{Reclaim Policy}


    \section{Understand Retain Reclaim Policy}


    \section{Storage Class}

%    \chapter{Exam Preparation Section}

\end{document}